\chapter*{Abstract}
\label{ch:abstract}

The growing generation of food waste and the surplus of crude glycerol from biodiesel production represent two significant organic waste streams requiring sustainable management solutions. This study evaluated the technical feasibility of a single-stage anaerobic co-digestion system treating food waste and glycerol as co-substrates, with the objective of verifing biogas production while providing a valorisation route for crude glycerol. The investigation was conducted in a 1 L continuous stirred-tank reactor operated under mesophilic conditions (35 °C) with a hydraulic retention time of 30 days. Food waste was collected from a university canteen, and crude glycerol was obtained from biodiesel production residues; anaerobic sludge from a municipal wastewater treatment plant served as inoculum. The reactor was operated in sequential phases: an initial adaptation period, co-digestion of food waste and glycerol, glycerol only feeding, and a post-eeding period. Biogas production was monitored continuously using an automatic methane potential test system, while process stability was assessed through measurements of pH, total alkalinity, volatile fatty acids, and their ratio. Substrate characterisation included total solids, volatile solids, density, total carbon, and total nitrogen analyses. Results demonstrated that co-digestion of food waste and glycerol yielded the highest methane production rates and cumulative volumes, with flow spikes immediately following feeding events indicating rapid substrate conversion. In contrast, glycerol-only feeding resulted in reduced biogas production despite continued methanogenic activity. Total alkalinity remained elevated following bicarbonate supplementation, maintaining pH around 7 throughout most of the experiment; volatile fatty acid concentrations fluctuated between 3000–5000 mg CH\textsubscript{3}COOH/L, with VFA to alkalinity ratios of 0.3-1.2. These findings confirm the hypothesis that glycerol can be useful for biogas production when co-digested with food waste or alone. The study demonstrates the technical viability of anaerobic co-digestion as a strategy for valorising crude glycerol within a circular bioeconomy framework, though further research is needed to optimise glycerol loading rates and develop robust feeding strategies for industrial application.
